% Created 2024-05-08 Wed 18:59
% Intended LaTeX compiler: pdflatex
\documentclass[presentation]{beamer}
\usepackage[utf8]{inputenc}
\usepackage[T1]{fontenc}
\usepackage{graphicx}
\usepackage{longtable}
\usepackage{wrapfig}
\usepackage{rotating}
\usepackage[normalem]{ulem}
\usepackage{amsmath}
\usepackage{amssymb}
\usepackage{capt-of}
\usepackage{hyperref}
\mode<beamer>{\usetheme{Madrid}}
\definecolor{SUred}{rgb}{0.59375, 0, 0.17969} % SU red (primary)
\definecolor{SUblue}{rgb}{0, 0.17578, 0.38281} % SU blue (secondary)
\setbeamercolor{palette primary}{bg=SUred,fg=white}
\setbeamercolor{palette secondary}{bg=SUblue,fg=white}
\setbeamercolor{palette tertiary}{bg=SUblue,fg=white}
\setbeamercolor{palette quaternary}{bg=SUblue,fg=white}
\setbeamercolor{structure}{fg=SUblue} % itemize, enumerate, etc
\setbeamercolor{section in toc}{fg=SUblue} % TOC sections
% Override palette coloring with secondary
\setbeamercolor{subsection in head/foot}{bg=SUblue,fg=white}
\setbeamercolor{date in head/foot}{bg=SUblue,fg=white}
\institute[SU]{Shenandoah University}
\titlegraphic{\includegraphics[width=0.5\textwidth]{\string~/Documents/suLogo/suLogo.pdf}}
\newcommand{\R}{\mathbb{R}}
\usepackage{tikz}
\usetheme{default}
\author{Chase Mathison\thanks{cmathiso@su.edu}}
\date{9 May 2024}
\title{Parametric Equations}
\hypersetup{
 pdfauthor={Chase Mathison},
 pdftitle={Parametric Equations},
 pdfkeywords={},
 pdfsubject={},
 pdfcreator={Emacs 29.1 (Org mode 9.6.7)}, 
 pdflang={English}}
\begin{document}

\maketitle

\section{Announcements}
\label{sec:org96cc938}
\begin{frame}[label={sec:org72d1bab}]{Announcements}
\begin{enumerate}
\item Quizz in Canvas.
\item Exam corrections due tomorrow.
\item Final Review in Canvas.
\item Final Exam, Wednesday at 8am.
\end{enumerate}
\end{frame}

\section{Lecture}
\label{sec:orgb36e4b8}
\begin{frame}[label={sec:org4431e78}]{Parametric equations}
How are we used to plotting equations?  

Usually, we are given an equation or a function, such as
\[
y = x^2 + 2x + 1 \]
or
\[
x^2 + y^2 = 4 \]
And we solve for \(y\) and plot using a table of values (if you don't
remember what the graph looks like).  

In other words, we usually have 1 \uline{\hspace*{1in}} variable and 1
\uline{\hspace*{1in}} variable.


But in real life, very often both \(x\) and \(y\) (and other
variables) depend on some other independent variable called a
\uline{\hspace*{1in}}.

This parameter is very often time, but it could be other things too.
\end{frame}

\begin{frame}[label={sec:org4e58c15}]{Example}
You go on a trip to run some errands.  Here is the data for your
position at different times.

\begin{center}
\begin{tabular}{|c|c|c|}
\hline
\(t\) & \(x\) & \(y\)\\
\hline
0 & 0 & 0\\
1 & 1 & -1\\
2 & 2 & -1\\
3 & 2 & 1\\
4 & 0 & 0\\
\hline
\end{tabular}
\end{center}

Assuming you traveled in straight lines between coordinates, draw a
plot of the path you took on your trip.

\vspace{10in}
\end{frame}

\begin{frame}[label={sec:org5c34e93}]{Example}
\begin{tikzpicture}
\draw[step = 0.5,black,thin] (0,0) grid (5,5);
\end{tikzpicture}
\end{frame}

\begin{frame}[label={sec:orge7a4b28}]{Parametric Equations}
It's a fact that any curve in the \(xy-\)plane can be represented
with \uline{\hspace*{1in}}.

\begin{definition}[Parametric curve]
If \(x\) and \(y\) are continuous functions of \(t\) on an interval
\(I\), then the equations
\[x = x \left( t \right), \quad y = y \left( t \right) \]
are called parametric equations and \(t\) is called the
\underline{\hspace{1in}}. The set of points \(\left( x,y \right)\)
obtained as \(t\) varies over the interval \(I\) is called the \underline{\hspace{1in}}
of the parametric equations.  The graph of parametric equations is
called a \underline{\hspace{1in}} or \emph{plane curve}, and is denoted by \(C\).
\end{definition}
\end{frame}

\begin{frame}[label={sec:org713bb88}]{Eliminating the Parameter}
Plotting parametric equations by making tables is a fine place to
start, but could give us misleading results.  For instance, if we try
to plot

\[
x \left( \theta \right) = \cos \left( 4\theta \right), \quad y \left(
\theta \right) = \sin \left( 4\theta \right), \quad 0 \le \theta \le 2\pi \]
And we use a table with multiples of \(\pi/4\), we might be mislead
into thinking that \(y\) is never changing.

For this reason, often the best way to plot a parametric curve is to
try to \uline{\hspace*{1in}} to end up with 1 equation involving
only \(x\) and \(y\).
\end{frame}

\begin{frame}[label={sec:org2fd592e}]{Example}
Let's eliminate the parameter \(t\) to plot the curve given by the
following parametric equations in the plane:
\[
x \left( t \right) = t^2, \quad y \left( t \right) = t - 4, \quad -1
\le t \le 2 \]
\vspace{10in}
\end{frame}

\begin{frame}[label={sec:org4cb9a91}]{Example}
\end{frame}

\begin{frame}[label={sec:orgfd35746}]{Example}
Eliminate the parameter \(\theta\) to plot the curve given by the
following parametric equations in the plane:
\[
x \left( \theta \right) = 8 \cos \left( \theta \right), \quad y \left(
\theta \right) = 8 \sin \left( \theta \right), \quad 0 \le \theta \le
2\pi \]
\vspace{10in}
\end{frame}

\begin{frame}[label={sec:org98a6550}]{Example}
\end{frame}

\begin{frame}[label={sec:org043ca4d}]{Example}
Let's see how you can develop parametric equations by picking a
parameter to work with using the example of a cycloid:
\vspace{10in}
\end{frame}

\begin{frame}[label={sec:orga7ed723}]{Example}
\end{frame}
\end{document}